Kurs:Algebraische Kurven (Osnabrück 2012)/Arbeitsblatt 25/latex
Kurs:Algebraische Kurven (Osnabrück 2012)/Arbeitsblatt 25/latex


\setcounter{section}{25}


  \zwischenueberschrift{Aufwärmaufgaben}


\inputaufgabegibtloesung

{}

{


Bestimme für die ebene algebraische Kurve


\mathdisp {V \left(  X^3+Y^2-XY+X  \right)} {  }


eine nichtkonstante Potenzreihenlösung


\mavergleichskette

{\vergleichskette

{ X
}

{ = }{ F(Y)
}

{  }{
}

{    }{
}

{   }{
}

}
{}{}{}
im Nullpunkt bis zum sechsten Glied.


}

{} {}


\inputaufgabegibtloesung

{}

{


Bestimme für die
\definitionsverweis {ebene algebraische Kurve}{}{}


\mathdisp {V \left(  X^2Y+X^2+Y^2-5XY+Y  \right)} {  }


eine nicht-konstante Potenzreihenlösung


\mavergleichskette

{\vergleichskette

{ Y
}

{ = }{ F(X)
}

{  }{
}

{    }{
}

{   }{
}

}
{}{}{}
im Nullpunkt bis zur fünften Ordnung.


}

{} {}


Die folgenden Aufgaben beschäftigen sich mit der Komplettierung eines lokalen Ringes.


\inputaufgabe

{}

{


Betrachte zu einem
\definitionsverweis {lokalen Ring}{}{}
$R$ mit
\definitionsverweis {maximalem Ideal}{}{}
${\mathfrak m}$ das Diagramm


\mathdisp {\longrightarrow R/{\mathfrak m}^4 \longrightarrow R/{\mathfrak m}^3 \longrightarrow R/{\mathfrak m}^2 \longrightarrow R/{\mathfrak m}} { . }


Dabei sind die Abbildungen die kanonischen Projektionen
\maabb {\varphi_{n}} { R/{\mathfrak m}^{n+1} } { R/{\mathfrak m}^n
} {,}
die durch die Idealinklusionen


\mavergleichskette

{\vergleichskette

{ {\mathfrak m}^{n+1}
}

{ \subseteq }{ {\mathfrak m}^n
}

{  }{
}

{    }{
}

{   }{
}

}
{}{}{}
induziert werden. Eine Folge von Elementen


\mavergleichskette

{\vergleichskette

{ a_n
}

{ \in }{ R/{\mathfrak m}^n
}

{  }{
}

{    }{
}

{   }{
}

}
{}{}{}
heißt \stichwort {verträglich} {,} wenn


\mavergleichskette

{\vergleichskette

{ \varphi_n(a_{n+1})
}

{ = }{ a_n
}

{  }{
}

{    }{
}

{   }{
}

}
{}{}{}
für alle $n$ gilt. Definiere eine Ringstruktur auf der Menge aller verträglichen Elemente
\zusatzklammer {diesen Ring nennt man die

\definitionswortenp{Komplettierung}{} von $R$} {} {.}
Zeige ferner, dass es einen kanonischen Ringhomomorphismus von $R$ in die Komplettierung gibt.


}

{} {}


\inputaufgabe

{}

{


Es sei $R$ ein eindimensionaler lokaler noetherscher kommutativer Ring. Zeige, dass die kanonische Abbildung von $R$ in die Komplettierung von $R$ injektiv ist.


}

{Bemerkung: Die Injektivität gilt für jeden noetherschen lokalen Ring, ist aber schwieriger zu beweisen.} {}


\inputaufgabe

{}

{


Es sei $R$ ein
\definitionsverweis {kommutativer Ring}{}{}
und $I$ ein
\definitionsverweis {Ideal}{}{.}
Zeige, dass durch


\mathdisp {\{x + I^n \, \vert \, n \in \N \} \quad (x \in R)} {  }


Umgebungsbasen definiert werden. Zeige außerdem, dass die auf $R$ induzierte Topologie genau dann hausdorffsch ist, wenn


\mavergleichskette

{\vergleichskette

{ \bigcap_n I^n
}

{ = }{ \{0\}
}

{  }{
}

{    }{
}

{   }{
}

}
{}{}{.}


}

{Bemerkung: Die Komplettierung eines lokalen Ringes bezüglich seines maximalen Ideals entspricht dann genau der (topologischen) Komplettierung bezüglich dieser Topologie.} {}


  \zwischenueberschrift{Aufgaben zum Abgeben}


\inputaufgabe

{}

{


Betrachte die Kardioide


\mathdisp {V \left(  \left(  X^2+Y^2  \right)^2- 2X \left(  X^2+Y^2  \right) -Y^2  \right)} {  }


im Punkt $(2,0)$. Bestimme eine formale Parametrisierung (bis zum fünften Term) der Kurve in diesem Punkt in Abhängigkeit von einem Tangentenparameter.


}

{} {}


\inputaufgabe

{}

{


Betrachte den Einheitskreis


\mavergleichskette

{\vergleichskette

{ X^2+Y^2
}

{ = }{ 1
}

{  }{
}

{    }{
}

{   }{
}

}
{}{}{}
im Punkt $(1,0)$. Bestimme Potenzreihen $G$ und


\mavergleichskette

{\vergleichskette

{ H
}

{ \in }{ K [ \![ T ]\! ]
}

{  }{
}

{    }{
}

{   }{
}

}
{}{}{}
mit den Anfangsbedingungen

\mathl{a_0=1,\, a_1=0,\, b_0=0,\, b_1=1}{} und mit


\mavergleichskette

{\vergleichskette

{ G(T)^2+H(T)^2
}

{ = }{ 1
}

{  }{
}

{    }{
}

{   }{
}

}
{}{}{.}


}

{} {}


\inputaufgabe

{}

{


Betrachte die Neilsche Parabel


\mavergleichskette

{\vergleichskette

{ C
}

{ = }{ V \left(  Y^3-X^2  \right)
}

{  }{
}

{    }{
}

{   }{
}

}
{}{}{}
im Punkt $(1,1)$. Finde eine Parametrisierung der Kurve in diesem Punkt mit Potenzreihen
\zusatzklammer {bis zum fünften Glied} {} {}
derart, dass eine Potenzreihe davon ein lineares Polynom ist.


}

{} {}


\inputaufgabe

{}

{


Es sei $K$ ein
\definitionsverweis {Körper}{}{.}
Eine

\definitionswortenp{formale Laurentreihe mit endlichem Hauptteil}{} ist eine unendliche Summe der Form


\mathdisp {F= \sum_{n=k}^\infty a_n T^n \text{ mit } a_n \in K \text{ und } k \in \Z} { . }


Zeige, dass der Ring dieser formalen Reihen
\zusatzklammer {mit geeigneten Ringoperationen} {} {}
isomorph zum
\definitionsverweis {Quotientenkörper}{}{}
des
\definitionsverweis {Potenzreihenringes}{}{}


\mathl{K [ \![ T ]\! ]}{} ist.


}

{} {}


\inputaufgabe

{}

{


Es sei $K$ ein
\definitionsverweis {Körper}{}{}
und sei $K[T]$ der
\definitionsverweis {Polynomring in einer Variablen}{}{.}
Es sei $R$ die
\definitionsverweis {Lokalisierung}{}{}
von $K[T]$ am
\definitionsverweis {maximalen Ideal}{}{}


\mavergleichskette

{\vergleichskette

{ {\mathfrak m}
}

{ = }{ (T)
}

{  }{
}

{    }{
}

{   }{
}

}
{}{}{.}
Zeige, dass die
\definitionsverweis {Komplettierung}{}{}
von $R$ isomorph zum
\definitionsverweis {Potenzreihenring}{}{}
$K [ \![ T]\! ]$ ist.


}

{} {}


\inputaufgabe

{}

{


Es sei $K$ ein
\definitionsverweis {Körper}{}{}
und


\mavergleichskette

{\vergleichskette

{ R
}

{ = }{ K [ \![ T ]\! ]
}

{  }{
}

{    }{
}

{   }{
}

}
{}{}{}
der
\definitionsverweis {Potenzreihenring}{}{}
Zeige, dass es in $R$ keine Quadratwurzel für $T$ gibt. Zeige ferner, dass für


\mavergleichskette

{\vergleichskette

{ K
}

{ = }{ \Z/( 7 )
}

{  }{
}

{    }{
}

{   }{
}

}
{}{}{}
das Element

\mathl{T+2}{} eine Quadratwurzel in $R$ besitzt, und bestimme die ersten fünf Koeffizienten von einer Quadratwurzel davon.


}

{} {}


\inputaufgabe

{}

{


Es sei


\mavergleichskette

{\vergleichskette

{ F
}

{ \in }{ K[X,Y]
}

{  }{
}

{    }{
}

{   }{
}

}
{}{}{}
ein irreduzibles Polynom und


\mavergleichskette

{\vergleichskette

{ R
}

{ = }{ K[X,Y]/(F)
}

{  }{
}

{    }{
}

{   }{
}

}
{}{}{}
der integre Koordinatenring der ebenen Kurve


\mavergleichskette

{\vergleichskette

{ C
}

{ = }{ V(F)
}

{  }{
}

{    }{
}

{   }{
}

}
{}{}{.}
Es sei

\mathl{R \rightarrow S = R ^{\operatorname{norm} }}{} die Normalisierung von $R$ und es sei

\mathl{R \rightarrow K [ \![ T]\! ]}{} der Ringhomomorphismus zu einer nichtkonstanten formalen Potenzreihenlösung der Kurve. Zeige, dass es einen eindeutig bestimmten Ringhomomorphismus

\mathl{S \rightarrow K [ \![ T]\! ]}{} gibt derart, dass das Diagramm


\mathdisp {\begin{matrix} R   & \stackrel{  }{\longrightarrow} &  S   & \\ &  \searrow & \downarrow & \\ & &  K [ \![ T]\! ]   & \!\!\!\!\! \!\!\!  \\ \end{matrix}} {  }


kommutiert.


}

{} {}


  \zwischenueberschrift{Aufgabe zum Hochladen}


\inputaufgabe

{}

{


Zeichne mittels eines geeigneten Programms eine der Beispielkurven der Vorlesung sowie die verschiedenen dort berechneten polynomialen Approximationen.


}

{} {}


<< | Kurs:Algebraische Kurven (Osnabrück 2012) | >>


PDF-Version dieses Arbeitsblattes


 Zur Vorlesung
(PDF)
